% !TeX program = xelatex
\ifdefined\BookMaster\else
\documentclass[11pt,oneside]{article}
% !TeX program = xelatex
\usepackage[b5paper,left=22mm,right=22mm,top=20mm,bottom=22mm,footskip=11mm]{geometry}
\usepackage{fontspec}
\setmainfont{TeX Gyre Termes}
\setsansfont{TeX Gyre Heros}
\usepackage{amsmath,amssymb,mathtools,amsthm}
\renewcommand{\cdots}{\mathinner{\cdotp\cdotp\cdotp}}
\usepackage{graphicx,adjustbox,array,booktabs,longtable,multirow,enumitem,changepage,needspace}
\usepackage[table]{xcolor}
\usepackage{fancyhdr}
\raggedbottom
\interfootnotelinepenalty=10000
\newcommand{\BookFootnote}[2]{\begingroup\renewcommand{\thefootnote}{\fnsymbol{footnote}}\footnote[#1]{#2}\endgroup}
\newcommand{\BookFootnoteMark}[1]{\begingroup\renewcommand{\thefootnote}{\fnsymbol{footnote}}\footnotemark[#1]\endgroup}
\newcommand{\BookFootnoteText}[2]{\begingroup\renewcommand{\thefootnote}{\fnsymbol{footnote}}\footnotetext[#1]{#2}\endgroup}
\newenvironment{BookFloat}[1]{\par\addvspace{.6\baselineskip}\noindent\def\BookFloatID{#1}\begin{minipage}{\linewidth}\phantomsection\label{bookfloat:#1:start}}{\phantomsection\label{bookfloat:\BookFloatID:end}\end{minipage}\par\addvspace{.6\baselineskip}}

\usepackage[unicode,hidelinks]{hyperref}
\pagestyle{fancy}
\fancyhf{}
\newcommand{\BookFooter}{}
\fancyfoot[L]{\parbox[b]{\dimexpr\textwidth-15mm}{\footnotesize\BookFooter}}
\fancyfoot[R]{\footnotesize\thepage}
\renewcommand{\headrulewidth}{0pt}
\renewcommand{\footrulewidth}{0pt}
\setlength{\parindent}{0pt}
\setlength{\parskip}{0.2em}
\setlength{\emergencystretch}{3em}
\setlist{nosep,leftmargin=1.7em,labelsep=.65em}
\newif\ifBookNarrative
\BookNarrativetrue
\newenvironment{BookBody}{\setlength{\parindent}{0pt}}{}
\newcommand{\BookStatement}[1]{\par\addvspace{.6\baselineskip}\noindent #1\enspace\ignorespaces}
\newcommand{\BookLabel}[1]{\par\addvspace{.6\baselineskip}\needspace{3\baselineskip}\noindent\textbf{#1}\quad\ignorespaces}
\newcommand{\BookHeading}[1]{\par\addvspace{\baselineskip}\needspace{4\baselineskip}\noindent{\large\bfseries #1}\par\medskip}
\newcommand{\BookSubheading}[1]{\par\addvspace{.7\baselineskip}\needspace{3\baselineskip}\noindent{\bfseries #1}\par\smallskip}
\newcommand{\BookTextHeading}[1]{\par\addvspace{.7\baselineskip}\needspace{3\baselineskip}\noindent #1\par\smallskip}
\newcommand{\BookSectionTitle}[2]{\par\addvspace{1.2\baselineskip}\needspace{5\baselineskip}\noindent\colorbox{gray!25}{\sffamily\bfseries\small SECTION #1}\quad{\large\bfseries #2}\par\medskip}
\newcommand{\BookExerciseTitle}[1]{\par\addvspace{\baselineskip}\needspace{4\baselineskip}\noindent{\color{gray}\rule{1.5mm}{1.5mm}}\enspace{\large\bfseries Exercises #1}\par\medskip}
\newcommand{\BookBullet}[1]{\begin{list}{\textbullet}{\setlength{\leftmargin}{1.7em}\setlength{\labelwidth}{1em}\setlength{\labelsep}{.7em}\setlength{\topsep}{0pt}\setlength{\parsep}{0pt}\setlength{\itemsep}{0pt}}\item #1\end{list}}
\newcommand{\BookExercise}[2]{\begin{list}{\textbf{#1}}{\setlength{\leftmargin}{2.2em}\setlength{\labelwidth}{1.5em}\setlength{\labelsep}{.7em}\setlength{\topsep}{.15em}\setlength{\parsep}{0pt}\setlength{\itemsep}{0pt}}\item #2\end{list}}
\newcommand{\BookEnd}[1]{\unskip\nobreak\hfill\ensuremath{#1}}
\newcommand{\BookSource}[2]{\phantomsection\label{sourcepage:#1}} % Source-page anchors, retained through reflow.
\usepackage{multicol}
\usepackage{tikz}
\newcommand{\BookPartEntry}[2]{\par\noindent\makebox[17mm][l]{\bfseries Section #1}#2\par}
\newcommand{\BookPartTitle}[3]{\par\bigskip\noindent\begin{tikzpicture}\node[anchor=north west,inner sep=0pt,outer sep=0pt] (bookpart) at (0,0) {\begin{minipage}{\dimexpr\linewidth-.5pt}\hspace*{26mm}\begin{minipage}[t]{\dimexpr\linewidth-26mm}\begin{adjustbox}{max width=\linewidth}{\fontsize{20}{24}\selectfont\bfseries #2}\end{adjustbox}\par\vspace{6mm}\hrule\vspace{6mm}{\sffamily\footnotesize\setlength{\parskip}{.3em}#3}\vspace{3mm}\end{minipage}\end{minipage}};\path[shade,top color=gray!60,bottom color=white] (0,0) rectangle ([xshift=24mm]bookpart.south west);\node[anchor=north,font=\sffamily\bfseries\small] at (12mm,-2mm) {PART};\node[anchor=north,font=\sffamily\bfseries\Huge] at (12mm,-7mm) {#1};\end{tikzpicture}\par\bigskip}
\newsavebox{\BookMathBox}
\makeatletter
\newcommand{\BookInline}[1]{\leavevmode\hskip0pt\relax\begingroup\in@{\begin}{#1}\ifin@\sbox{\BookMathBox}{\(#1\)}\ifdim\wd\BookMathBox>\dimexpr\linewidth-1em\relax\resizebox{\dimexpr\linewidth-1em\relax}{!}{\usebox{\BookMathBox}}\else\(#1\)\fi\else\in@{\left}{#1}\ifin@\sbox{\BookMathBox}{\(#1\)}\ifdim\wd\BookMathBox>\dimexpr\linewidth-1em\relax\resizebox{\dimexpr\linewidth-1em\relax}{!}{\usebox{\BookMathBox}}\else\(#1\)\fi\else\(#1\)\fi\fi\endgroup\hskip0pt\relax}
\makeatother
\DeclareMathOperator{\Aut}{Aut}
\DeclareMathOperator{\End}{End}
\DeclareMathOperator{\Hom}{Hom}
\DeclareMathOperator{\GF}{GF}
\DeclareMathOperator{\irr}{irr}
\DeclareMathOperator{\lt}{lt}
\DeclareMathOperator{\lp}{lp}
\providecommand{\textcircled}[1]{\textnormal{(#1)}}

\begin{document}
\fi
\renewcommand{\BookFooter}{Section 21 Group Presentations}
\pdfbookmark[0]{Section 21 Group Presentations}{21_group_presentations}
\ifdefined\BookMaster\else\setcounter{page}{177}\fi
% Source PDF page 194; original printed page 177
\BookSource{194}{177}
\BookNarrativetrue

\begin{BookBody}

\BookSectionTitle{21}{Group Presentations}

\BookTextHeading{\textbf{Definition}}

\noindent Following most of the literature on group presentations, in this section we let 1 be the identity of a group. The idea of a\textit{ group presentation} is to form a group by giving a set of generators for the group and certain equations or relations that we want the generators to satisfy. We want the group to be as free as it possibly can be on the generators, subject to these relations.

\end{BookBody}

% Source PDF page 195; original printed page 178
\BookSource{195}{178}
\BookNarrativetrue

\begin{BookBody}

\BookLabel{21.1 Example} Suppose \BookInline{G} is\textit{ free with generation} \BookInline{x} \textit{and} \BookInline{y} \textit{except for the relation} \BookInline{xy= yx,\allowbreak } which we may express as \BookInline{xyx^{- 1}y^{- 1}= 1.} Note that the condition \BookInline{xy= yx} is exactly what is needed to make \BookInline{G} abelian, even though \BookInline{xyx^{- 1}y^{- 1}} is just one of the many possible commutators of \BookInline{F[\{x,\allowbreak y\}].} Thus \BookInline{G} is free abelian on two generators and is isomorphic to \BookInline{F[\{x,\allowbreak y\}]} modulo its commutator subgroup. This commutator subgroup of \BookInline{F[\{x,\allowbreak y\}]} is the smallest normal subgroup containing \BookInline{xyx^{- 1}y^{- 1},\allowbreak } since any normal subgroup containing \BookInline{xyx^{- 1}y^{- 1}} gives rise to a factor group that is abelian and thus contains the commutator subgroup by Theorem \BookInline{13.22.}\BookEnd{\blacktriangle}

The preceding example illustrates the general situation. Let \BookInline{F[A]} be a free group and suppose that we want to form a new group as much like \BookInline{F[A]} as it can be, subject to certain equations that we want satisfied. Any equation can be written in a form in which the right-hand side is 1. Thus we can consider the equations to be \BookInline{r_{i}= 1} for \BookInline{i\in I,\allowbreak } where \BookInline{r_{i}\in F[A].} If we require that \BookInline{r_{i}= 1,\allowbreak } then we will have to have

\begin{center}\begin{adjustbox}{max width=\linewidth}\BookInline{\displaystyle x\Big(r_{i}^{n}\Big)x^{- 1}= 1}\end{adjustbox}\end{center}

\noindent for any \BookInline{x\in F[A]} and \BookInline{n\in \mathbb{Z}.} Also any product of elements equal to 1 will again have to equal 1. Thus any finite product of the form

\begin{center}\begin{adjustbox}{max width=\linewidth}\BookInline{\displaystyle \prod _{j}x_{j}\Big(r_{i_{j}}^{n_{j}}\Big)x_{j}^{- 1},\allowbreak }\end{adjustbox}\end{center}

\noindent where the \BookInline{r_{i_{j}}} need not be distinct, will have to equal 1 in the new group. It is readily checked that the set of all these finite products is a normal subgroup \BookInline{R} of \BookInline{F[A].} Thus any group looking as much as possible like \BookInline{F[A],\allowbreak } subject to the requirements \BookInline{r_{i}= 1,\allowbreak } also has \BookInline{r= 1} for every \BookInline{r\in R.} But \BookInline{F[A]/R} looks like \BookInline{F[A](}remember that we multiply cosets by choosing representatives), except that \BookInline{R} has been collapsed to form the identity 1. Hence the group we are after is (at least isomorphic to\BookInline{)\; F[A]/R.} We can view this group as described by the generating set \BookInline{A} and the set \BookInline{\{r_{i}\mid i\in I\},\allowbreak } which we will abbreviate \BookInline{\{r_{i}\}.}

\BookTextHeading{\ensuremath{\blacksquare} \textbf{HISTORICAL NOTE}  }

\noindent The idea of a group presentation already appears in Arthur Cayley’s 1859 paper, “On the Theory of Groups as Depending on the Symbolic Equation \BookInline{\theta ^{n}= 1.} Third Part.” In this article, Cayley gives a complete enumeration of the five groups of order 8, both by listing all the elements of each and by giving for each a presentation. For example, his third example is what is here called the\textit{ dihedral group} \BookInline{D_{4};} Cayley notes that this group is generated by the two elements \BookInline{\alpha,\allowbreak \beta } with the relations \BookInline{\alpha ^{4}= 1,\allowbreak \beta ^{2}= 1,\allowbreak \alpha \beta = \beta \alpha ^{3}.} He also shows more generally that a group of order \BookInline{mn} is generated by \BookInline{\alpha,\allowbreak \beta } with the relations \BookInline{\alpha ^{m}= 1,\allowbreak } \BookInline{\beta ^{n}= 1,\allowbreak \alpha \beta = \beta \alpha ^{s}} if and only if \BookInline{s^{n}\equiv 1(}mod \BookInline{m)} (see Exercise 13).

\noindent In 1878, Cayley returned to the theory of groups and noted that a central problem in that theory is the determination of all groups of a given order \BookInline{n.} In the early 1890s, Otto Ho¨lder published several papers attempting to solve Cayley’s problem. Using techniques similar to those discussed in Sections 17 and 21, Ho¨lder determined all simple groups of order up to 200 and characterized all the groups of orders \BookInline{p^{3},\allowbreak pq^{2},\allowbreak pqr,\allowbreak } and \BookInline{p^{4},\allowbreak } where \BookInline{p,\allowbreak q,\allowbreak r} are distinct prime numbers. Furthermore, he developed techniques for determining the possible structures of a group \BookInline{G,\allowbreak } if one is given the structure of a normal subgroup \BookInline{H} and the structure of the factor group \BookInline{G/H.} Interestingly, since the notion of an abstract group was still fairly new at this time, Ho¨lder typically began his papers with the definition of a group and also emphasized that isomorphic groups are essentially one and the same object.

\BookLabel{21.2 Definition} Let \BookInline{A} be a set and let \BookInline{\{r_{i}\}\subseteq F[A].} Let \BookInline{R} be the least normal subgroup of \BookInline{F[A]} containing the \BookInline{r_{i}.} An isomorphism \BookInline{\phi } of \BookInline{F[A]/R} onto a group \BookInline{G} is a\textbf{ presentation of} \BookInline{G.} The 
% Source PDF page 196; original printed page 179
\BookSource{196}{179}
sets \BookInline{A} and \BookInline{\{r_{i}\}} give a\textbf{ group presentation.} The set \BookInline{A} is the set of\textbf{ generators for the} \textbf{presentation} and each \BookInline{r_{i}} is a\textbf{ relator.} Each \BookInline{r\in R} is a\textbf{ consequence of} \BookInline{\{r_{i}\}.} An equation \BookInline{r_{i}= 1} is a\textbf{ relation.} A\textbf{ finite presentation} is one in which both \BookInline{A} and \BookInline{\{r_{i}\}} are finite sets\BookInline{.}\BookEnd{\blacksquare}

\end{BookBody}

\BookNarrativetrue

\begin{BookBody}

This definition may seem complicated, but it really is not. In Example \BookInline{21.1,\allowbreak \{x,\allowbreak y\}} is our set of generators and \BookInline{xyx^{- 1}y^{- 1}} is the only relator. The equation \BookInline{xyx^{- 1}y^{- 1}= 1,\allowbreak } or \BookInline{xy= yx,\allowbreak } is a relation. This was an example of a finite presentation.

If a group presentation has generators \BookInline{x_{j}} and relators \BookInline{r_{i},\allowbreak } we shall use the notations

\begin{center}\begin{adjustbox}{max width=\linewidth}\BookInline{\displaystyle (x_{j}:r_{i})\;\text{or}\;(x_{j}:r_{i}= 1)}\end{adjustbox}\end{center}

\noindent to denote the group presentation. We may refer to \BookInline{F[\{x_{j}\}]/R} as\textit{ the group with presenta}\textit{tion} \BookInline{(x_{j}:r_{i}).}

\BookTextHeading{\textbf{Isomorphic Presentations}  }

\BookLabel{21.3 Example} Consider the group presentation with

\begin{center}\begin{adjustbox}{max width=\linewidth}\BookInline{\displaystyle A= \{a\}\;\text{and}\;\{r_{i}\}= \{a^{6}\},\allowbreak }\end{adjustbox}\end{center}

\noindent that is, the presentation

\begin{center}\begin{adjustbox}{max width=\linewidth}\BookInline{\displaystyle (a:a^{6}= 1).}\end{adjustbox}\end{center}

\noindent This group defined by one generator \BookInline{a,\allowbreak } with the relation \BookInline{a^{6}= 1,\allowbreak } is isomorphic to \BookInline{\mathbb{Z}_{6}.}

Now consider the group defined by two generators \BookInline{a} and \BookInline{b,\allowbreak } with \BookInline{a^{2}= 1,\allowbreak b^{3}= 1,\allowbreak } and \BookInline{ab= ba,\allowbreak } that is, the group with presentation

\begin{center}\begin{adjustbox}{max width=\linewidth}\BookInline{\displaystyle (a,\allowbreak b:a^{2},\allowbreak b^{3},\allowbreak aba^{- 1}b^{- 1}).}\end{adjustbox}\end{center}

\noindent The condition \BookInline{a^{2}= 1} gives \BookInline{a^{- 1}= a.} Also \BookInline{b^{3}= 1} gives \BookInline{b^{- 1}= b^{2}.} Thus every element in this group can be written as a product of nonnegative powers of \BookInline{a} and \BookInline{b.} The relation \BookInline{aba} \BookInline{^ { - 1 } b ^ { - 1 } = 1}, that is\BookInline{,\allowbreak ab= ba,\allowbreak } allows us to write first all the factors involving \BookInline{a} and then the factors involving \BookInline{b.} Hence every element of the group is equal to some \BookInline{a^{m}b^{n}.} But then \BookInline{a^{2}= 1} and \BookInline{b^{3}= 1} show that there are just six distinct elements,

\begin{center}\begin{adjustbox}{max width=\linewidth}\BookInline{\displaystyle 1,\allowbreak b,\allowbreak b^{2},\allowbreak a,\allowbreak ab,\allowbreak ab^{2}.}\end{adjustbox}\end{center}

\noindent The subgroup \BookInline{\langle ab\rangle } contains the elements \BookInline{1,\allowbreak ab,\allowbreak } and the powers of \BookInline{ab:}

\begin{center}\begin{adjustbox}{max width=\linewidth}\BookInline{\displaystyle \begin{aligned}(ab)^{2}&= a^{2}b^{2}= b^{2} \\ (ab)^{3}&= abb^{2}= a \\ (ab)^{4}&= a(ab)= b \\ (ab)^{5}&= (ab)b= ab^{2}.\end{aligned}}\end{adjustbox}\end{center}

\noindent So this group is also a cyclic group of order 6 isomorphic with \BookInline{\mathbb{Z}_{6}.} \BookEnd{\blacktriangle}

\BookLabel{21.4 Example} The dihedral group has presentation

\begin{center}\begin{adjustbox}{max width=\linewidth}\BookInline{\displaystyle D_{n}:(a,\allowbreak b|a^{n},\allowbreak b^{2},\allowbreak abab)}\end{adjustbox}\end{center}

\noindent since if we let \BookInline{a= \mu } and \BookInline{b= \rho } the three relations are exactly the defining relations for \BookInline{D_{n}.(}The last relation \BookInline{abab= 1} is equivalent to \BookInline{ab= ba^{- 1}.)} The element \BookInline{abab} is in \BookInline{R} if and only if \BookInline{b(abab)b^{- 1}} is in \BookInline{R} since \BookInline{R} is a normal subgroup. We have \BookInline{b(abab)b^{- 1}= } \BookInline{baba.} So in any presentation with generators \BookInline{a} and \BookInline{b} and a relator \BookInline{abab,\allowbreak } we can replace \BookInline{abab} with \BookInline{baba} and get the same subgroup \BookInline{R} and therefore the same factor group. Hence the dihedral group also has presentation

\begin{center}\begin{adjustbox}{max width=\linewidth}\BookInline{\displaystyle D_{n}:(a,\allowbreak b:a^{n},\allowbreak b^{2},\allowbreak baba).}\end{adjustbox}\end{center}

\end{BookBody}

% Source PDF page 197; original printed page 180
\BookSource{197}{180}
\BookNarrativetrue

\begin{BookBody}

\noindent Setting the relators to 1 gives the equivalent presentation

\begin{center}\begin{adjustbox}{max width=\linewidth}\BookInline{\displaystyle (a,\allowbreak b:a^{n}= 1,\allowbreak b^{2}= 1,\allowbreak ba= a^{- 1}b^{- 1})}\end{adjustbox}\end{center}

\noindent which can also be rewritten as

\begin{center}\begin{adjustbox}{max width=\linewidth}\BookInline{\displaystyle (a,\allowbreak b:a^{n}= 1,\allowbreak b^{2}= 1,\allowbreak ba= a^{n- 1}b).}\end{adjustbox}\end{center} \BookEnd{\blacktriangle}

The preceding examples illustrate that different presentations may give isomorphic groups. When this happens, we have\textbf{ isomorphic presentations.} To determine whether two presentations are isomorphic may be very hard. It has been shown (see Rabin [22]) that a number of such problems connected with this theory are not generally solvable; that is, there is no\textit{ routine} and well-defined way of discovering a solution in all cases. These unsolvable problems include the problem of deciding whether two presentations are isomorphic, whether a group given by a presentation is finite, free, abelian, or trivial, and the famous\textit{ word problem} of determining whether a given word \BookInline{w} is a consequence of a given set of relations \BookInline{\{r_{i}\}.}

The importance of this material is indicated by our Theorem 20.13, which guarantees that\textit{ every group has} \BookInline{a} \textit{presentation}.

\BookLabel{21.5 Example} Let us show that

\begin{center}\begin{adjustbox}{max width=\linewidth}\BookInline{\displaystyle (x,\allowbreak y:y^{2}x= y,\allowbreak yx^{2}y= x)}\end{adjustbox}\end{center}

\noindent is a presentation of the trivial group of one element. We need only show that \BookInline{x} and \BookInline{y} are consequences of the relators \BookInline{y^{2}xy^{- 1}} and \BookInline{yx^{2}yx^{- 1},\allowbreak } or that \BookInline{x= 1} and \BookInline{y= 1} can be deduced from \BookInline{y^{2}x= y} and \BookInline{yx^{2}y= x.} We illustrate both techniques.

As a consequence of \BookInline{y^{2}xy^{- 1},\allowbreak } we get \BookInline{yx} upon conjugation by \BookInline{y^{- 1}.} From \BookInline{yx} we deduce \BookInline{x^{- 1}y^{- 1},\allowbreak } and then \BookInline{(x^{- 1}y^{- 1})(yx^{2}yx^{- 1})} gives \BookInline{xyx^{- 1}.} Conjugating \BookInline{xyx^{- 1}} by \BookInline{x^{- 1},\allowbreak } we get \BookInline{y.} From \BookInline{y} we get \BookInline{y^{- 1},\allowbreak } and \BookInline{y^{- 1}(yx)} is \BookInline{x.}

Working with relations instead of relators, from \BookInline{y^{2}x= y} we deduce \BookInline{yx= 1} upon multiplication by \BookInline{y^{- 1}} on the left. Then substituting \BookInline{yx= 1} into \BookInline{yx^{2}y= x,\allowbreak } that is, \BookInline{(yx)(xy)= x,\allowbreak } we get \BookInline{xy= x.} Then multiplying by \BookInline{x^{- 1}} on the left, we have \BookInline{y= 1.} Substituting this in \BookInline{yx= 1,\allowbreak } we get \BookInline{x= 1.}

Both techniques amount to the same work, but it somehow seems more natural to most of us to work with relations\BookInline{.}\BookEnd{\blacktriangle}

\BookTextHeading{\textbf{Applications}}

\noindent We conclude this chapter with two applications.

\BookLabel{21.6 Example} Let us determine all groups of order 10 up to isomorphism. We know from the Fundamental Theorem 9.12 that every abelian group of order 10 is isomorphic to \BookInline{\mathbb{Z}_{10}.} Suppose that \BookInline{G} is nonabelian of order 10. By Sylow theory\BookInline{,\allowbreak G} contains a normal subgroup \BookInline{H} of order 5, and \BookInline{H} must be cyclic. Let \BookInline{a} be a generator of \BookInline{H.} Then \BookInline{G/H} is of order 2 and thus isomorphic to \BookInline{Z_{2}.} If \BookInline{b\in G} and \BookInline{b\notin H,\allowbreak } we must then have \BookInline{b^{2}\in H.} Since every element of \BookInline{H} except 1 has order 5, if \BookInline{b^{2}} were not equal to 1, then \BookInline{b^{2}} would have order 5, so \BookInline{b} would have order 10. This would mean that \BookInline{G} would be cyclic, contradicting our assumption that \BookInline{G} is not abelian. Thus \BookInline{b^{2}= 1.} Finally, since \BookInline{H} is a normal subgroup of \BookInline{G,\allowbreak bHb^{- 1}= H,\allowbreak } so in particular\BookInline{,\allowbreak bab^{- 1}\in H.} Since conjugation by \BookInline{b} is an automorphism of \BookInline{H,\allowbreak bab^{- 1}} must be another element of \BookInline{H} of order 5, hence \BookInline{bab^{- 1}} equals \BookInline{a,\allowbreak a^{2},\allowbreak a^{3},\allowbreak } or \BookInline{a^{4}.} But \BookInline{bab^{- 1}= a} would give \BookInline{ba= ab,\allowbreak } and then \BookInline{G} would be abelian, since \BookInline{a} and \BookInline{b} generate \BookInline{G.} Thus the possibilities for presentations of \BookInline{G} are:

\BookExercise{1.}{\BookInline{(a,\allowbreak b:a^{5}= 1,\allowbreak b^{2}= 1,\allowbreak ba= a^{2}b),\allowbreak }}

\BookExercise{2.}{\BookInline{(a,\allowbreak b:a^{5}= 1,\allowbreak b^{2}= 1,\allowbreak ba= a^{3}b),\allowbreak }}

\BookExercise{3.}{\BookInline{(a,\allowbreak b:a^{5}= 1,\allowbreak b^{2}= 1,\allowbreak ba= a^{4}b).}}

\end{BookBody}

% Source PDF page 198; original printed page 181
\BookSource{198}{181}
\BookNarrativetrue

\begin{BookBody}

Note that all three of these presentations can give groups of order at most 10, since the last relation \BookInline{ba= a^{i}b} enables us to express every product of \BookInline{a}’s and \BookInline{b}’s in \BookInline{G} in the form \BookInline{a^{s}b^{t}.} Then \BookInline{a^{5}= 1} and \BookInline{b^{2}= 1} show that the set

\begin{center}\begin{adjustbox}{max width=\linewidth}\BookInline{\displaystyle S= \{a^{0}b^{0},\allowbreak a^{1}b^{0},\allowbreak a^{2}b^{0},\allowbreak a^{3}b^{0},\allowbreak a^{4}b^{0},\allowbreak a^{0}b^{1},\allowbreak a^{1}b^{1},\allowbreak a^{2}b^{1},\allowbreak a^{3}b^{1},\allowbreak a^{4}b^{1}\}}\end{adjustbox}\end{center}

\noindent includes all elements of \BookInline{G.}

It is not yet clear that all these elements in \BookInline{S} are distinct, so that we have in all three cases a group of order 10. For example, the group presentation

\begin{center}\begin{adjustbox}{max width=\linewidth}\BookInline{\displaystyle (a,\allowbreak b:a^{5}= 1,\allowbreak b^{2}= 1,\allowbreak ba= a^{2}b)}\end{adjustbox}\end{center}

\noindent gives a group in which, using the associative law, we have

\begin{center}\begin{adjustbox}{max width=\linewidth}\BookInline{\displaystyle \begin{aligned}a&= b^{2}a= (bb)a= b(ba)= b(a^{2}b)= (ba)(ab) \\ &= (a^{2}b)(ab)= a^{2}(ba)b= a^{2}(a^{2}b)b= a^{4}b^{2}= a^{4}\end{aligned}}\end{adjustbox}\end{center}

\noindent Thus in this group\BookInline{,\allowbreak a= a^{4},\allowbreak } so \BookInline{a^{3}= 1,\allowbreak } which, together with \BookInline{a^{5}= 1,\allowbreak } yields \BookInline{a^{2}= 1.} But \BookInline{a^{2}= 1,\allowbreak } together with \BookInline{a^{3}= 1,\allowbreak } means that \BookInline{a= 1.} Hence every element in the group with presentation

\begin{center}\begin{adjustbox}{max width=\linewidth}\BookInline{\displaystyle (a,\allowbreak b:a^{5}= 1,\allowbreak b^{2}= 1,\allowbreak ba= a^{2}b)}\end{adjustbox}\end{center}

\noindent is equal to either 1 or \BookInline{b;} that is, this group is isomorphic to \BookInline{\mathbb{Z}_{2}.} A similar study of

\begin{center}\begin{adjustbox}{max width=\linewidth}\BookInline{\displaystyle (bb)a= b(ba)}\end{adjustbox}\end{center}

\noindent for

\begin{center}\begin{adjustbox}{max width=\linewidth}\BookInline{\displaystyle (a,\allowbreak b:a^{5}= 1,\allowbreak b^{2}= 1,\allowbreak ba= a^{3}b)}\end{adjustbox}\end{center}

\noindent shows that \BookInline{a= a^{4}} again, so this also yields a group isomorphic to \BookInline{\mathbb{Z}_{2}.}

This leaves just

\begin{center}\begin{adjustbox}{max width=\linewidth}\BookInline{\displaystyle (a,\allowbreak b:a^{5}= 1,\allowbreak b^{2}= 1,\allowbreak ba= a^{4}b)}\end{adjustbox}\end{center}

\noindent as a candidate for a nonabelian group of order 10. As in Example 21.4 this is a presentation of the dihedral group \BookInline{D_{5}.}

If we were unaware of the dihedral group, how would we show that the presentation gives a group with 10 elements? One attack is as follows. Let us try to make \BookInline{S} into a group by defining \BookInline{(a^{s}b^{t})(a^{u}b^{v})} to be \BookInline{a^{x}b^{y},\allowbreak } where \BookInline{x} is the remainder of \BookInline{s+ u(4^{t})} when divided by 5, and \BookInline{y} is the remainder of \BookInline{t+ v} when divided by 2, in the sense of the division algorithm (Theorem 6.2). The formula \BookInline{s+ u(4)^{t}} is counting what the power of \BookInline{a} should be after moving \BookInline{u} copies of \BookInline{a} by \BookInline{t} copies of \BookInline{b.} In other words, we use the relation \BookInline{ba= a^{4}b} as a guide in defining the product \BookInline{(a^{s}b^{t})(a^{u}b^{v})} of two elements of \BookInline{S.} We see that \BookInline{a^{0}b^{0}} acts as identity, and that given \BookInline{a^{u}b^{v},\allowbreak } we can determine \BookInline{t} and \BookInline{s} successively by letting

\begin{center}\begin{adjustbox}{max width=\linewidth}\BookInline{\displaystyle t\equiv - v(\;\text{mod}\;2)}\end{adjustbox}\end{center}

\noindent and then

\begin{center}\begin{adjustbox}{max width=\linewidth}\BookInline{\displaystyle s\equiv - u(4^{t})(\;\text{mod}\;5),\allowbreak }\end{adjustbox}\end{center}

\noindent giving \BookInline{a^{s}b^{t},\allowbreak } which is a left inverse for \BookInline{a^{u}b^{v}.} We will then have a group structure on \BookInline{S} if and only if the associative law holds. Exercise 13 asks us to carry out the straightforward computation for the associative law and to discover a condition for \BookInline{S} to be a group under such a definition of multiplication. The criterion of the exercise in this case amounts to the valid congruence

\begin{center}\begin{adjustbox}{max width=\linewidth}\BookInline{\displaystyle 4^{2}\equiv 1(\;\text{mod}\;5).}\end{adjustbox}\end{center}

\noindent Thus we do get a group of order 10. Note that

\begin{center}\begin{adjustbox}{max width=\linewidth}\BookInline{\displaystyle 2 ^ { 2 } \not \equiv 1 \left( \bmod 5 \right)}\end{adjustbox}\end{center}

\end{BookBody}

% Source PDF page 199; original printed page 182
\BookSource{199}{182}
\BookNarrativetrue

\begin{BookBody}

\noindent and

\begin{center}\begin{adjustbox}{max width=\linewidth}\BookInline{\displaystyle 3 ^ { 2 } \not \equiv 1 \pmod { 5 } ,\allowbreak }\end{adjustbox}\end{center}

\noindent so Exercise 13 also shows that

\begin{center}\begin{adjustbox}{max width=\linewidth}\BookInline{\displaystyle (a,\allowbreak b:a^{5}= 1,\allowbreak b^{2}= 1,\allowbreak ba= a^{2}b)}\end{adjustbox}\end{center}

\noindent and

\begin{center}\begin{adjustbox}{max width=\linewidth}\BookInline{\displaystyle (a,\allowbreak b:a^{5}= 1,\allowbreak b^{2}= 1,\allowbreak ba= a^{3}b)}\end{adjustbox}\end{center}

\noindent do not give groups of order 10. \BookEnd{\blacktriangle}

\BookLabel{21.7 Example} Let us determine all groups of order 8 up to isomorphism. We know the three abelian ones:

\begin{center}\begin{adjustbox}{max width=\linewidth}\BookInline{\displaystyle \mathbb{Z}_{8},\allowbreak \mathbb{Z}_{2}\times \mathbb{Z}_{4},\allowbreak \mathbb{Z}_{2}\times \mathbb{Z}_{2}\times \mathbb{Z}_{2}.}\end{adjustbox}\end{center}

\noindent Using generators and relations, we shall give presentations of the nonabelian groups.

Let \BookInline{G} be nonabelian of order 8. Since \BookInline{G} is nonabelian, it has no elements of order 8, so each element but the identity is of order either 2 or 4. If every element were of order 2, then for \BookInline{a,\allowbreak b\in G,\allowbreak } we would have \BookInline{(ab)^{2}= 1,\allowbreak } that is\BookInline{,\allowbreak abab= 1.} Then since \BookInline{a^{2}= 1} and \BookInline{b^{2}= 1} also, we would have

\begin{center}\begin{adjustbox}{max width=\linewidth}\BookInline{\displaystyle ba= a^{2}bab^{2}= a(ab)^{2}b= ab,\allowbreak }\end{adjustbox}\end{center}

\noindent contrary to our assumption that \BookInline{G} is not abelian. Thus \BookInline{G} must have an element of order 4.

Let \BookInline{\langle a\rangle } be a subgroup of \BookInline{G} of order 4. If \BookInline{b\notin \langle a\rangle,\allowbreak } the cosets \BookInline{\langle a\rangle } and \BookInline{b\langle a\rangle } exhaust all of \BookInline{G.} Hence \BookInline{a} and \BookInline{b} are generators for \BookInline{G} and \BookInline{a^{4}= 1.} Since \BookInline{\langle a\rangle } is normal in \BookInline{G(}by Sylow theory, or because it is of index \BookInline{2),\allowbreak G/\langle a\rangle } is isomorphic to \BookInline{\mathbb{Z}_{2}} and we have \BookInline{b^{2}\in \langle a\rangle.} If \BookInline{b^{2}= a} or \BookInline{b^{2}= a^{3},\allowbreak } then \BookInline{b} would be of order 8. Hence \BookInline{b^{2}= 1} or \BookInline{b^{2}= a^{2}.} Finally, since \BookInline{\langle a\rangle } is normal, we have \BookInline{bab^{- 1}\in \langle a\rangle,\allowbreak } and since \BookInline{b\langle a\rangle b^{- 1}} is a subgroup conjugate to \BookInline{\langle a\rangle } and hence isomorphic to \BookInline{\langle a\rangle,\allowbreak } we see that \BookInline{bab^{- 1}} must be an element of order 4. Thus \BookInline{bab^{- 1}= a} or \BookInline{bab^{- 1}= a^{3}.} If \BookInline{bab^{- 1}} were equal to \BookInline{a,\allowbreak } then \BookInline{ba} would equal \BookInline{ab,\allowbreak } which would make \BookInline{G} abelian. Hence \BookInline{bab^{- 1}= a^{3},\allowbreak } so \BookInline{ba= a^{3}b.} Thus we have two possibilities for \BookInline{G,\allowbreak } namely,

\begin{center}\begin{adjustbox}{max width=\linewidth}\BookInline{\displaystyle G_{1}:(a,\allowbreak b:a^{4}= 1,\allowbreak b^{2}= 1,\allowbreak ba= a^{3}b)}\end{adjustbox}\end{center}

\noindent and

\begin{center}\begin{adjustbox}{max width=\linewidth}\BookInline{\displaystyle G_{2}:(a,\allowbreak b:a^{4}= 1,\allowbreak b^{2}= a^{2},\allowbreak ba= a^{3}b).}\end{adjustbox}\end{center}

Note that \BookInline{a^{- 1}= a^{3},\allowbreak } and that \BookInline{b^{- 1}} is \BookInline{b} in \BookInline{G_{1}} and \BookInline{b^{3}} in \BookInline{G_{2}.} These facts, along with the relation \BookInline{ba= a^{3}b,\allowbreak } enable us to express every element in \BookInline{G_{i}} in the form \BookInline{a^{m}b^{n},\allowbreak } as in Examples 21.3 and 21.6. Since \BookInline{a^{4}= 1} and either \BookInline{b^{2}= 1} or \BookInline{b^{2}= a^{2},\allowbreak } the possible elements in each group are

\begin{center}\begin{adjustbox}{max width=\linewidth}\BookInline{\displaystyle 1,\allowbreak a,\allowbreak a^{2},\allowbreak a^{3},\allowbreak b,\allowbreak ab,\allowbreak a^{2}b,\allowbreak a^{3}b.}\end{adjustbox}\end{center}

\noindent Thus \BookInline{G_{1}} and \BookInline{G_{2}} each have order at most 8. The first group \BookInline{G_{1}} is sometimes called the \textbf{octic} group, but as we saw in Example 21.4 it is isomorphic with our old friend \BookInline{D_{4},\allowbreak } the dihedral group. For the second we can make \BookInline{S} into a group by defining \BookInline{(a^{i}b^{j})(a^{r}b^{s})} to be \BookInline{a^{x}b^{y}} where \BookInline{y} is \BookInline{j+ s} modulo 2 and if \BookInline{j+ s<2,\allowbreak } then \BookInline{x} is the remainder of \BookInline{i+ r(2j+ 1)} when divided by 4 and if \BookInline{j+ s= 2,\allowbreak } then \BookInline{x} is the remainder when \BookInline{i+ 2+ r(2j+ 1)} is divided by 4. We leave it as an exercise to show that this operation makes \BookInline{S} a group, which shows that \BookInline{G_{2}} is a presentation of a group of order 8.

Since \BookInline{ba= a^{3}b\neq ab,\allowbreak } we see that both \BookInline{G_{1}} and \BookInline{G_{2}} are nonabelian. That the two groups are not isomorphic follows from the fact that a computation shows that \BookInline{G_{1}} has only two elements of order 4, namely\BookInline{,\allowbreak a} and \BookInline{a^{3}.} On the other hand, in \BookInline{G_{2}} all elements but 1 and \BookInline{a^{2}} are of order 4. We leave the computations of the tables for these groups 
% Source PDF page 200; original printed page 183
\BookSource{200}{183}
to Exercise 3. To illustrate suppose we wish to compute \BookInline{(a^{2}b)(a^{3}b).} Using \BookInline{ba= a^{3}b} repeatedly, we get

\end{BookBody}

\BookNarrativetrue

\begin{BookBody}

\begin{center}\begin{adjustbox}{max width=\linewidth}\BookInline{\displaystyle (a^{2}b)(a^{3}b)= a^{2}(ba)a^{2}b= a^{5}(ba)ab= a^{8}(ba)b= a^{11}b^{2}.}\end{adjustbox}\end{center}

\noindent Then for \BookInline{G_{1},\allowbreak } we have

\begin{center}\begin{adjustbox}{max width=\linewidth}\BookInline{\displaystyle a^{11}b^{2}= a^{11}= a^{3},\allowbreak }\end{adjustbox}\end{center}

\noindent but if we are in \BookInline{G_{2},\allowbreak } we get

\begin{center}\begin{adjustbox}{max width=\linewidth}\BookInline{\displaystyle a^{11}b^{2}= a^{13}= a.}\end{adjustbox}\end{center}

The group \BookInline{G_{2}} is called the\textbf{ quaternion group}. We shall encounter the quaternion group again in Section \BookInline{32.}\BookEnd{\blacktriangle}

\end{BookBody}

\BookNarrativefalse

\begin{BookBody}

\BookExerciseTitle{21}

\BookTextHeading{\textbf{Computations}}

\BookExercise{1.}{Give a presentation of \BookInline{\mathbb{Z}_{4}} involving one generator; involving two generators; involving three generators.}

\BookExercise{2.}{Give a presentation of \BookInline{S_{3}} involving three generators.}

\BookExercise{3.}{Give the tables for both the octic group}

\begin{center}\begin{adjustbox}{max width=\linewidth}\BookInline{\displaystyle (a,\allowbreak b:a^{4}= 1,\allowbreak b^{2}= 1,\allowbreak ba= a^{3}b)}\end{adjustbox}\end{center}

\noindent and the quaternion group

\begin{center}\begin{adjustbox}{max width=\linewidth}\BookInline{\displaystyle (a,\allowbreak b:a^{4}= 1,\allowbreak b^{2}= a^{2},\allowbreak ba= a^{3}b).}\end{adjustbox}\end{center}

\noindent In both cases, write the elements in the order \BookInline{1,\allowbreak a,\allowbreak a^{2},\allowbreak a^{3},\allowbreak b,\allowbreak ab,\allowbreak a^{2}b,\allowbreak a^{3}b.(}Note that we do not have to compute\textit{ every} product. We know that these presentations give groups of order 8, and once we have computed enough products the rest are forced so that each row and each column of the table has each element exactly once.)

\BookExercise{4.}{Determine all groups of order 14 up to isomorphism. [\textit{Hint:} Follow the outline of Example 21.6 and use Exercise 13, part (b).]}

\BookExercise{5.}{Determine all groups of order 21 up to isomorphism. [\textit{Hint:} Follow the outline of Example 21.6 and use Exercise 13, part (b). It may seem that there are two presentations giving nonabelian groups. Show that they are isomorphic.]}

\BookTextHeading{\textbf{Concepts}}

\noindent In Exercises 6 and 7, correct the definition of the italicized term without reference to the text, if correction is needed, so that it is in a form acceptable for publication.

\BookExercise{6.}{A\textit{ consequence} of the set of relators is any finite product of relators raised to powers.}

\BookExercise{7.}{Two group presentations are\textit{ isomorphic} if and only if there is a one-to-one correspondence of the generators of the first presentation with the generators of the second that yields, by renaming generators, a one-to-one correspondence of the relators of the first presentation with those of the second.}

\BookExercise{8.}{Determine whether each of the following is true or false.}

\BookExercise{a.}{Every group has a presentation.}

\BookExercise{b.}{Every group has many different presentations.}

\BookExercise{c.}{Every group has two presentations that are not isomorphic.}

\BookExercise{d.}{Every group has a finite presentation.}

\BookExercise{e.}{Every group with a finite presentation is of finite order.}

\BookExercise{f.}{Every cyclic group has a presentation with just one generator.}

\BookExercise{g.}{Every conjugate of a relator is a consequence of the relator.}

\end{BookBody}

% Source PDF page 201; original printed page 184
\BookSource{201}{184}
\BookNarrativefalse

\begin{BookBody}

\BookExercise{h.}{Two presentations with the same number of generators are always isomorphic.}

\BookExercise{i.}{In a presentation of an abelian group, the set of consequences of the relators contains the commutator subgroup of the free group on the generators.}

\BookExercise{j.}{Every presentation of a free group has 1 as the only relator.}

\BookTextHeading{\textbf{Theory}}

\BookExercise{9.}{Use the methods of this section and Exercise 13, part (b), to show that there are no nonabelian groups of order 15.}

\BookExercise{10.}{Show, using Exercise 13, that}

\begin{center}\begin{adjustbox}{max width=\linewidth}\BookInline{\displaystyle (a,\allowbreak b:a^{3}= 1,\allowbreak b^{2}= 1,\allowbreak ba= a^{2}b)}\end{adjustbox}\end{center}

\noindent gives a group of order 6. Show that it is nonabelian.

\BookExercise{11.}{Show that the presentation}

\begin{center}\begin{adjustbox}{max width=\linewidth}\BookInline{\displaystyle (a,\allowbreak b:a^{3}= 1,\allowbreak b^{2}= 1,\allowbreak ba= a^{2}b)}\end{adjustbox}\end{center}

\noindent of Exercise 10 gives (up to isomorphism) the only nonabelian group of order 6, and hence gives a group isomorphic to \BookInline{S_{3}.}

\BookExercise{12.}{We showed in Example 13.6 that \BookInline{A_{4}} has no subgroup of order 6. The preceding exercise shows that such a subgroup of \BookInline{A_{4}} would have to be isomorphic to either \BookInline{\mathbb{Z}_{6}} or \BookInline{S_{3}.} Show again that this is impossible by considering orders of elements.}

\BookExercise{13.}{Let}

\begin{center}\begin{adjustbox}{max width=\linewidth}\BookInline{\displaystyle S= \{a^{i}b^{j}\mid 0\leq i<m,\allowbreak 0\leq j<n\},\allowbreak }\end{adjustbox}\end{center}

\noindent that is\BookInline{,\allowbreak S} consists of all formal products \BookInline{a^{i}b^{j}} starting with \BookInline{a^{0}b^{0}} and ending with \BookInline{a^{m- 1}b^{n- 1}.} Let \BookInline{r} be a positive integer, and define multiplication on \BookInline{S} by

\begin{center}\begin{adjustbox}{max width=\linewidth}\BookInline{\displaystyle (a^{s}b^{t})(a^{u}b^{v})= a^{x}b^{y},\allowbreak }\end{adjustbox}\end{center}

\noindent where \BookInline{x} is the remainder of \BookInline{s+ u(r^{t})} when divided by \BookInline{m,\allowbreak } and \BookInline{y} is the remainder of \BookInline{t+ v} when divided by \BookInline{n,\allowbreak } in the sense of the division algorithm (Theorem 6.2).

\BookExercise{a.}{Show that a necessary and sufficient condition for the associative law to hold and for \BookInline{S} to be a group under this multiplication is that \BookInline{r^{n}\equiv 1(}mod \BookInline{m).}}

\BookExercise{b.}{Deduce from part (a) that the group presentation}

\begin{center}\begin{adjustbox}{max width=\linewidth}\BookInline{\displaystyle (a,\allowbreak b:a^{m}= 1,\allowbreak b^{n}= 1,\allowbreak ba= a^{r}b)}\end{adjustbox}\end{center}

gives a group of order \BookInline{mn} if and only if \BookInline{r^{n}\equiv 1(}mod \BookInline{m).(}See the Historical Note in this section.)

\BookExercise{14.}{Without using Exercise 13, prove that \BookInline{(a,\allowbreak b:a^{5}= 1,\allowbreak b^{2}= 1,\allowbreak ba= a^{3}b)} is a presentation for the group \BookInline{\mathbb{Z}_{2}.}}

\BookExercise{15.}{Is the group obtained from the group presentation with the letters a through z as generators and the words in a standard English dictionary as relators the trivial group? Prove your answer.}

\end{BookBody}
\ifdefined\BookMaster\let\BookFinish\relax\else\def\BookFinish{\end{document}}\fi\BookFinish
